\chapter{Travail associé}
Cette partie présente ce qui se fait déjà dans le monde à l'heure actuelle en matière d'apprentissage de la programmation.

Premièrement, il y aura un tour d'horizon des politiques gouvernementales de pays qui ont adopté l'apprentissage de la programmation dans leur programme officiel. Ceci présentera les avancées de cette problématique à travers plusieurs pays et aidera à envisager différentes manières de la traiter. 

Après ce tour d'horizon, l'attention sera portée sur des organisations non gouvernementales qui ont pour but de promouvoir l'enseignement de la programmation. 

La troisième partie décrira les langages utilisés dans les organisations présentées précédemment. 

Enfin, des concepts différenciateurs seront extraits afin de distinguer la spécificité de chaque initiative.
\subsection{L'état de la programmation dans le monde}
\label{monde}
Cette partie fait un tour d'horizon de différent pays qui enseigne la programmation aux jeunes. Spécialement sur le positionnement politique que ces pays ont adopté par rapport à notre problématique. Cette partie ne porte donc plus sur une méthode ou une organisation, mais bien sûr le contexte et la manière d'intégré l'apprentissage de la programmation dans chaque pays.
\subsubsection{Angleterre}
L'apprentissage de l'informatique en Angleterre\footnote{\url{https://www.gov.uk/government/collections/statutory-guidance-schools\#national-curriculum-from-september-2014}} n'est pas nouveau. Cependant, longtemps cet apprentissage était centré sur les technologies de l'information et de la communication (TIC). En 2010, une étude a été commandée à la Royal Society pour évaluer cet apprentissage. Un an plus tard leur rapport a révélé que l'enseignement tel que dispensé jusque-là, n'était ni efficace ni en adéquation avec l'évolution de l'informatique dans notre société. La Royal Society suggère de changer les matières abordées en informatique. En effet précédemment c'était les TIC qui étaient enseignées alors que l'apprentissage de la programmation serait plus bénéfique et adapter pour les enfants. Sur base de ce rapport, les programmes de cours ont été adaptés.

\subsubsection{France}
En France\footnote{\url{http://fr.wikipedia.org/wiki/Informatique\_et\_sciences\_du\_num\%C3\%A9rique}
\url{http://fr.wikipedia.org/wiki/Baccalaur\%C3\%A9at\_scientifique}} depuis deux ans l'informatique fait partie intégrante du programme du baccalauréat de type S. Une des matières dispensées est "Informatique et sciences du numérique". Cette matière se subdivise en quatre sous parties qui sont : représentation de l'information, algorithmique, langage et programmation, architectures matérielles. Cette approche est donc également basée sur l'apprentissage de la programmation plutôt que sur les TIC.

\subsubsection{Nouvelle-Zélande}
Ce pays a adopté récemment les sciences informatiques dans son programme d'étude. Les cours sont dispensés à partir de 15 ans. Les cours dispensés concernent l'apprentissage de la programmation et des concepts informatique en général. Nous avons un choix un peu différent dans ce pays sur l'âge du début de l'apprentissage. Beaucoup de pays commencent plus jeunes et introduisent des concepts basiques. La Nouvelle-Zélande s'inscrit dans une logique plus similaire à la France, mais ne restreint pas l'informatique aux options scientifiques.

\subsubsection{Corée du Sud}
Enseigne l'informatique depuis longtemps et à tous les niveaux de l'enseignement. La culture numérique dans ces pays y est fort différente que par chez nous. Par exemple, une carrière dans le gaming y est tout à fait normale. L'informatique est vraiment omniprésente dans ces cultures, il est donc normal que son apprentissage commence en primaire.

\subsubsection{Grèce}
L'apprentissage des sciences informatiques prend une place importante dans les programmes grecs. Des 6 ans, les enfants sont confrontés à l'informatique à l'école. À cet âge, c'est plus de la maîtrise de l'outil qu'ils apprennent. Dès 10 ans, leurs cours d'informatique prennent une tournure plus algorithmique et donc plus proche de la science de l'informatique.

\section{Organisations existantes}
Il sera présenté ici un rassemblement d'initiatives similaires à ce travail. Ces initiatives sont toutes non gouvernementaux et propose d'enseigner la programmation de différentes manières. C'est sur ce dernier point que l'accent sera mis dans la suite de ce tour des organisations.
\subsection{Code.org}
\begin{figure}[!h]
  \begin{center}
    \includegraphics[scale=0.5]{content/5-related_work/images/code}
    \caption{Logo de code.org}
    \label{fig:code.org}
  \end{center}
\end{figure}
Code.org\footnote{\url{http://code.org/about}} est une organisation sans but lucratif des USA qui à pour objectifs :
\begin{itemize}
  \item apporter l'informatique dans toutes les classes de secondaire aux USA ;
  \item démontrer le succès de l'utilisation de cours en ligne dans l'enseignement public ;
  \item ajouter l'informatique dans les bases des programmes de sciences et de mathématique des 50 états ;
  \item employé la connaissance technique collective pour améliorer l'apprentissage de l'informatique dans le monde ;
  \item augmenter la représentation féminine et des personnes de couleurs dans informatique.
\end{itemize}

Pour ce faire, ils fournissent une plate-forme\footnote{\url{http://code.org/educate/20hr}} web qui permet aux professeurs de suivre leurs élèves grâce à un système de classe.

Toutes les ressources sont gratuites et librement utilisables\footnote{\url{http://code.org/faq}}. Leur programme d'apprentissage se base sur Blockly (voir \ref{blockly}).
Les ressources sont conçues pour que les professeurs comme les étudiants puissent commencer le cours sans connaître l'informatique (une assistance est proposée au professeur gratuitement si nécessaire).

Le site propose aux professeurs de se faire récompenser s'ils arrivent à finir les 27 missions proposées à minimum 15 étudiants. Dans ce cas, ils gagnent $750\$$, s'ils ont au moins 7 filles dans le groupe ils peuvent prétendre à $250\$$ supplémentaire.

\subsubsection{Déroulement des leçons}
Il propose des sessions de une heure de travail/jeu/apprentissage. Chaque unité d'une heure est découpée en petite mission (ex:5-20) les missions sont très courte et apporte un concept de programmation. Avant l'introduction de chaque concept, une petite vidéo est faite pour expliquer le concept introduit et des exemples de ce que la programmation permet de réaliser avec ce dernier.

Ils proposent de faire travailler les étudiants par pair\footnote{pair programming \url{http://en.wikipedia.org/wiki/Pair\_programming}}. Ceci permet d'avoir moins de questions pour le professeur et de mieux s'approprier la matière. Le travail par pair permet également de casser l'image du "geek" en montrant que la programmation est une science sociale et collaborative. Sans oublier qu'avec deux enfants par station, moins d'ordinateurs seront nécessaires.

Le site explique également que pour faire participer tout les élèves, il faut avoir confiance en leur compétence : permettre aux premiers groupes d'aider les derniers.

Quand un étudiant a une question, ils recommandent de proposer aux étudiants de d'abord demander à 3 de leurs camarades avant de poser la question au professeur. Le professeur ne devant pas être compétent, il doit juste pouvoir réfléchir avec les élèves de quel est le problème, cela permet aussi d'évité les questions de distraction ou de manque de compréhension.

Pour chaque petite mission, il y a un test automatisé qui informe si la mission est réussie ou non. Si la mission est réussie, le programme passe à la mission suivante. Il y a également un compteur de blocs dans les premières missions. Ce compteur permet de voir combien de blocs sont nécessaires pour réaliser la mission de manière optimale.

\subsection{CoderDojo}
\begin{figure}[!h]
  \begin{center}
    \includegraphics[scale=0.5]{content/5-related_work/images/dojo}
    \caption{Logo de CoderDojo}
    \label{fig:coder dojo}
  \end{center}
\end{figure}
CoderDojo\footnote{\url{http://coderdojo.com/about}} est un réseau open source de clubs de programmation dans le sens le plus large du terme. Tous les dojos sont donc autonomes. Dans ceux-ci, des enfants de 5 à 17 ans apprennent la programmation (site web, application, jeux...). La seul règle est "Above All : Be Cool" qui peut être mis en pratique simplement en créant des espaces d'échanges de savoir amical et sociable.

CoderDojo a été créé par James Whelton un irlandais de 18 ans et Bill Liao un entrepreneur australien à Cork. James a eu des demandes de jeunes enfants pour avoir des cours de programmation après qu'il eut hacké l'iPod nano. Beaucoup de gens de Dublin vinrent à ses cours et donc un nouveau Dojo a été créé à Dublin et puis cela s'est étendu à tout le globe.

\subsection{Code Club}
\begin{figure}[!h]
  \begin{center}
    \includegraphics[scale=0.3]{content/5-related_work/images/club}
    \caption{Logo de Code Club}
    \label{fig:code club}
  \end{center}
\end{figure}
Code Club\footnote{\url{https://www.codeclub.org.uk/about}} est un réseau de clubs national mené par des bénévoles en dehors des heures de cours. Ces activités s'adressent à des enfants de 9 à 11 ans.

Ils créent donc le matériel pour permettre à des bénévoles de donner des cours parascolaires d'environ une heure semaine. Il propose dans l'ordre d'utiliser scratch, HTML/css et puis Python. Ils aimeraient que les 21000 écoles primaires anglaises aient leur club.

Leur philosophie est de d'abord l'amusement, la créativité et l'exploration avant l'apprentissage des concepts de programmation.
\subsection{KidsCode}
\begin{figure}[!h]
  \begin{center}
    \includegraphics[scale=0.5]{content/5-related_work/images/kidscode}
    \caption{Logo de code.org}
    \label{fig:code.org}
  \end{center}
\end{figure}
KidsCode est une petite initiative wallonne qui fournit un stage de pour les enfants de 10 à 14 ans. Les ateliers sont dispensés par deux informaticiens qui proposent de partager leur passion. Ils coachent les jeunes de manière ludique dans le but de les rendre autonomes.

KidsCode propose d'apprendre la programmation grâce au langage Python. Ce langage a été choisi pour sa facilité d'utilisation et aussi du fait de son utilisation en condition réelle. 

KidsCode a été conçu et pensé en association avec l'incubateur de startup Nestup. Cette initiative est le fruit du constat, par le monde de l'entreprise, d'un manque d'intérêt pour la programmation de la part des plus jeunes. Ceci conduisant à un manque de vocation dans les startups.

\subsection{Langages de programmation}
Cette partie traite des différents éditeurs et langages de programmation utilisés dans les organisations présentées au point précédent.
\subsubsection{Blockly}
\label{blockly}
Blocky\footnote{\url{https://code.google.com/p/blockly/}} est un langage de programmation graphique basé sur des technologie du web. 

Blocky est influencé\footnote{\url{https://code.google.com/p/blockly/wiki/Alternatives}} par "App Inventor" qui est influencé par "Scratch". Ce dernier est lui influencé par "StarLogo".

Il a comme particularité:
\begin{enumerate}
\item De s'exécuter dans un navigateur;
\item D'exporter du code source en JavaScript, dart, etc..;
\item D'être open source;
\item D'être haut-niveau.
\end{enumerate}

Il n'est pas directement une plat-forme d'éducation dans le sens ou suivant les blocs implémenté, il sera utilisé pour l'éducation, le business, des jeux, etc..

Lors de la conception du langage de blocky, il devait avoir certaines propriétés. Les trois première sont pour augmenter la compréhension des néophytes, les autres portaient sur des facilités souhaitée du langage. Les propriétés décidées lors de la conception du langage sont\footnote{\url{https://code.google.com/p/blockly/wiki/Language}} :

\begin{itemize}
  \item Des index de liste commençant à 1;
  \item Des noms de variables non sensible à la casse;
  \item Pas de portée de variable, toutes les variables sont globale;
  \item La possibilité de pouvoir faire un export en JavaScript;
  \item Un code natif généré proche de celui des blocs.
\end{itemize}

\subsubsection{Scratch 1}
Scratch est un langage de programmation développé par le MIT pour apprendre aux enfants la programmation. C'est l'interface qui permet de faire des scripte facilement grâce à sa programmation par bloc.\\

Ce langage à été pensé pour être un outils créatif pour réaliser des histoires, des jeux, des simulations, de l'art, etc.. Ce langage a par exemple son propre éditeurs d'image et de sons. Un autre but est que ce langage doit être simple à utiliser et à apprendre. Il a en effet été conçu pour des enfants n'ayant aucune connaissance préalable en programmation.\\

L'interface de Scratch se divise en plusieurs grandes parties:
\begin{enumerate}
\item Sur la gauche il y a la zone de dessin dans laquelle s'anime les composants graphique des scriptes;
\item Au milieu il y a une liste des blocs disponible trier par catégorie.
\item Sur la droite il y a la zone de scripte qui contient tout les scriptes liés au sprit sélectionné.
\end{enumerate}

\subsubsection{SNAP}
SNAP est un langage de programmation de glissé-déposé de blocs. c'est une ré-implémentation et une extension du langage Scratch du MIT. C'est un langage qui a été pensé et conçu pour être orienté web, ce langage est donc implémenté en JavaScript.\\

Ce langage est né en 2011 et a été créer par Jens Mönig. Il se distingue de son père Scratch par l'ajout de :
\begin{enumerate}
\item Fonctions et des procédures de première classe;
\item Listes de première classe;
\item Sprite de première classe.
\end{enumerate}

\subsubsection{Python}
Python est un langage de programmation qu'il ne faut plus présenté. Il a beaucoup d'avantage dont celui d'avoir une syntaxe légère et d'être facile à prendre en main. Ceci couplé à une grande communauté et beaucoup de librairies dont la fameuse \texttt{turtle} en fait un excellent langage pour démarrer dans la programmation.

Le pendant d'utiliser un vrai langage de programmation est bien sûr de devoir l'apprendre en plus d'apprendre la logique de programmation. De plus il faut également enseigner les bonnes pratiques de codage la première fois que l'on demande à quelqu'un de programmer avec un vrai langage. 

C'est donc un choix qui ne convient pas aux trop jeunes ni à ceux qui ne sont pas vraiment investis dans l'apprentissage de la programmation.

\subsection{Concepts différenciateurs}
\label{concepts}
Maintenant à travers le fonctionnement des  différentes structures à travers le monde vu au point précédent, ce chapitre va mettre en lumière les différences et similitudes des méthodes d'apprentissage de la programmation. Ceci grâce à une série de critères tel que le publique cible, le lieu, les moyens et les personnes.
\subsubsection{pour qui ?}
L'apprentissage de la programmation aux enfants ne se fait pas de la même manière avec les différentes tranches d'âge. Les différentes organisations auront donc soit un public visé restreint soit différents programmes pour les différents groupes d'enfants. 

\paragraph{Age}
La méthode d'apprentissages doit être différents suivant l'âge. Les enfants n'ont pas toujours la maturité d'apprendre des concepts parfois fort abstrait (boucle, parallélisation...). En effet pour des élèves de maternel par exemple un tableau interactif sera un grand plus car ils n'ont pas encore toute la dextérité souhaitée sur la souris. Et au contraire des enfants plus grands seront plus autonomes et pourront travailler par paire voir seul.

\paragraph{Origine et genre}
Une autre différence est l'origine des enfants, certaines organisations, tel Code.org, mettent l'accent sur la disponibilité pour tous et donc notamment pour les filles qui sont sous représentée en informatique. Tout comme il est important de ne pas viser que les personnes qui peuvent acquérir un ordinateur à la maison ou se payer un stage. Dans ce cas, pour toucher un maximum de personne, l'école reste l'endroit le plus adéquat.

\subsubsection{Commment ?}
La principale différences entre toutes les organisations va être sur la manière d'aborder l'apprentissage. Les différences peuvent être de plusieurs ordres, tel que : les outils utilisés, les concepts abordés et l'environnement de travail.

\paragraph{Outils} 
Un aspect très pratique qui différencie l'apprentissage est les outils utilisés. Ici encore différentes écoles s'affrontent :
\begin{description}
  \item[Langage réel] Les enfants doivent apprendre ce qu'est la programmation et donc ce qu'est du code source. De plus, avec un langage déjà couramment utilisé, plus de choses sont possibles.
  \item[Langage visuel] Dans la programmation, le plus important est la logique. Donc pourquoi ne pas avoir un langage visuel par blocs pour simplifier la syntaxe et se rapprocher un peu de ce qui se fait généralement avec des diagrammes.
  \item[Langage web] Dans le monde d'aujourd'hui, le web est partout, plus que tout le reste. l'accent doit donc être porté sur l'apprentissage de ce que les enfants voient au quotidien.
\end{description}

\paragraph{Concepts} 
L'ordre dans laquelle les concepts de programmation sont abordés peut se différencier sur plusieurs points. Tous les cours se divisent en deux morceaux : la théorie et la pratique. Mais il existe de multiples manières d'assembler les deux dans un cours unique.
\begin{itemize}
  \item Certains commencent par donner de la théorie et demandent ensuite de l'appliquer ;
  \item D'autres proposent des exercices avant d'expliquer ce que les enfants ont utilisé de manière intuitive ;
  \item D'autres encore découpent la matière en tout petit sous ensemble avant d'appliquer la première ou la seconde technique.\\
  \end{itemize}


Ces différentes façons d'aborder l'apprentissage mènent à différentes façon de concevoir les missions. Celles-ci peuvent être très dense ou au contraire un succession de mission simple. Les missions dense sont souvent constituées de nombreux concepts de programmation (boucle, condition, fonction...) pour avoir directement un objectif concret alors que les missions simple essayent de bien séparé l'apprentissage de ceux-ci pour fournir une étude la plus progressive possible.\\

Dans certains pays, la programmation est associé au cours de sciences. Elle deviendra donc souvent un outil scientifique et perdra du même coup un peu de son coté ludique. Néanmoins, grâce à cela un plus grand nombre d'enfants pourront toucher à la programmation.\\

Dans tout les cas, il faudra trouver une manière ludique d'apprendre la programmation aux enfants. En effet, les enfants ont besoin de s'approprier la matière. Quoi de mieux que un jeux pour les garder attentif à ce qu'ils réalisent.

\paragraph{Environement de travail}
L'organisation d'une leçon peut varier dans la gestion du groupe :
\begin{description}
  \item[En groupe] L'apprentissage de groupe est souvent réalisé avec un tableau interactif ou un projecteur \footnote{\url{http://scratched.media.mit.edu/resources/scratch-\%C3\%A0-la-maternelle}} ;
  \item[Par binôme] La programmation en binôme est la plus répandue dans les cours de programmation. Faire travailler les enfants par deux, les obliges à interagir entre eux et à essayer de résoudre eux même les problèmes qu'ils rencontrent. cela permet aussi de diminuer le nombre d'ordinateurs utilisés ;
  \item[Individuel] Cette technique est très peu utilisé car elle demande beaucoup de moyen matériel et personnel ;
  \item[A domicile] De nombreuses plat-formes proposent des cours en lignes qui peuvent être suivi au rythme de chacun. Néanmoins, pour avoir de l'aide, il faudra se tourner vers une personne compétente ou vers les forums. Ces formations sont souvent plus théorique ou demande plus de lecture, ce qui empêche les jeunes enfants de les suivre.
\end{description}

\subsubsection{Par qui ?}
Les personnes impliquée peuvent être de différent types suivant les organisations. Ceux-ci peuvent être des enseignant formé ou non, des passionnés ou encore des professionnels. Au delà de des personnes il y a également le type de cours et leur création qui seront au centre de cette partie.

\paragraph{Type d'organisation}
Différentes organisations proposent des cours de programmations.On peut déjà différencier le fait que l'informatique fasse partie du programme officiel ou pas ou plus simplement que l'apprentissage se déroule à durant les cours ou en parascolaire.\\

Dans le premier cas, les étudiants ne viennent pas tous par choix contrairement au second. Ceci peut entraîner des blocage et des réticences à participer correctement à l'activité, comme cela peut arriver pour tout les autres cours dans une école.

Pour les cours donnés en parascolaire, aussi bien de petites structures très locales, se partagent le terrain avec de grandes organisations à visée internationale. Les problèmes typique de ce genre d'organisation est le coût des activités et les horaires.

\paragraph{Enseignants}
La plupart des cours sont donnés par des informaticiens ou au minimum des personnes formées à l'informatique. C'est bien entendu la manière la plus logique de travailler pour fournir un enseignement de qualité. Mais comme le nombre d'informaticien est réduit, certaines organisations propose des cours clés en mains pour des enseignants sans connaissances particulières en programmation. Ces professeurs apprendront en même temps que leur étudiants et seront la uniquement pour les aider à réfléchir.\\

Il n'y a pas de tendance globales sur le type de rémunération des enseignants. Ils recevront parfois des prime en fonction de la réussite de leur classe, d'autre seront payé directement par les étudiants ou par l'état, d'autres encore seront bénévoles.

\paragraph{Création des cours}
Suivant l'organisation, les cours seront plus ou moins figé. Certaines organisations permettent à leur communauté de modifier, créer des cours tandis que d'autres ont des équipes dédié à cette tache.


